\section{A.4. The letter to Bers}
\begin{proposition}[Refuted, with rechecked extremal-length certificate: the orange slices]\label{A4-orange-extremal}
Let $Q_0=\{|z|\leq1,\ \operatorname{Re}z\leq1/2\}$ in the sphere with centres $0,1,\infty$. The canonical cap $D_0$ for a generalized pant of boundary lengths $(32,0,0)$ is not contained in $Q_0$. In particular the literal model in the letter to Bers, which permits the other two caps to degenerate to points, is false.
\end{proposition}
\begin{proof}
We recheck the extremal-length argument credited in the input to Michel Hua's project. Suppose $D_0\subset Q_0$. For $z\in Q_0$,
\[
 |z+1/2|^2=|z|^2+\operatorname{Re}z+1/4\leq7/4.
\]
Since $\sqrt7/2<4/3<7/5<3/2$, the annulus
\[
 A=\{4/3<|z+1/2|<7/5\}
\]
lies in $P=\mathbf P^1\setminus(D_0^\circ\cup\{1,\infty\})$, and its core is freely homotopic in $P$ to the boundary of $D_0$. For this curve family $\gamma$, monotonicity under inclusion of the annulus gives
\[
\operatorname{Ext}_P(\gamma)\leq\operatorname{Ext}_A(\gamma)
=\frac{2\pi}{\log(21/20)}.
\]
The hyperbolic metric on $P$ has area $2\pi$ by Gauss--Bonnet, and the shortest curve in the boundary class has length $32$. In the defining supremum of extremal length it therefore gives
\[
\operatorname{Ext}_P(\gamma)\geq\frac{32^2}{2\pi}.
\]
These bounds contradict $1024\log(21/20)>4\pi^2$. For an exact elementary verification, $\log(21/20)=\int_{20}^{21}dt/t>1/21$, while $1024/21>4(22/7)^2>4\pi^2$. The annulus extremal length is obtained by the logarithmic coordinate, which gives a Euclidean rectangle with width $\log(21/20)$ and periodic height $2\pi$. Cauchy--Schwarz gives the matching upper bound for every conformal density. The two-cusp pant exists for every positive remaining boundary length. Every geometric and numerical hypothesis in the contradiction is thus satisfied.
\end{proof}

\begin{corollary}[Refuted: fixed concentric orange-slice caps]\label{A4-orange-jets}
There are no three fixed pointed discs $Q_0,Q_1,Q_\infty$ whose independently concentric subdiscs describe all canonically capped pants. Thus the literal model in 67/019--020 is false, independently of its additional spherical description.
\end{corollary}
\begin{proof}
If $U:Q_0\to\D$ is a centred uniformization, the unit-disc coordinate on the concentric subdisc $U^{-1}(\lambda\D)$ is $U/\lambda$. Its coefficient $b$ in Proposition \ref{A2-schwarzian} is independent of $\lambda$. The two actual pants in Theorem \ref{A2-independence} have different $b$.
\end{proof}

\subsection{The canonical bundle furnished by grafting}
\begin{theorem}[Known with derivation: capping is infinite grafting]\label{A4-grafting}
Let $S$ be a marked compact oriented surface of type $(g,n)$ with $n\geq1$ and negative Euler characteristic. Let $\mathcal T_B^{\geq0}(S)$ classify complete hyperbolic structures with each labelled end either a geodesic boundary of length $\ell_i>0$ or a cusp of length $\ell_i=0$, and with no nodes in the interior. Let $\mathcal T_{g,n}$ be its punctured Teichm\"uller space. Canonical capping gives a mapping-class-equivariant homeomorphism
\[
\Phi=(\operatorname{cap},\ell):\mathcal T_B^{\geq0}(S)
\longrightarrow\mathcal T_{g,n}\times[0,\infty)^n.
\]
Its restriction to each boundary-type stratum is a real-analytic diffeomorphism. In particular the positive-length part is a trivial principal $(\R_{>0})^n$-bundle over $\mathcal T_{g,n}$ by boundary lengths, not by all possible tangent radii.
\end{theorem}
\begin{proof}
The conformal cylinder $(\R/\ell_i\Z)\times[0,\infty)$ is the punctured unit disc through $u=\exp(2\pi i(s+iy)/\ell_i)$. Adding its missing point gives exactly the pointed disc attached by normalized arclength in 67/065--066. Thus the cap operation is Mondello's $\operatorname{gr}_\infty$, not just an analogy with it. Apply Mondello, Section 5.4, Theorem 5.4 \cite{AMondello2011}, restricted to the finite, non-nodal strata. His length vector is the same $\ell$ and his marking convention is the one obtained by collapsing boundary circles to the punctures. At zero total length his blown-up decoration directions are forgotten; the continuous maps descend to the displayed product. For all positive lengths the theorem supplies both the inverse and real analyticity, and on each mixed zero/positive stratum it supplies the same assertion. Hence the comparison transfers exactly the claimed structure. A section is $Y\mapsto\Phi^{-1}(Y,(1,\ldots,1))$.
\end{proof}

\begin{corollary}[Proved: an intrinsic action and retraction]\label{A4-action}
The multiplicative monoid $[0,\infty)^n$ acts continuously and intrinsically on $\mathcal T_B^{\geq0}(S)$ by
\[
\lambda\cdot X=\Phi^{-1}\bigl(\operatorname{cap}(X),(\lambda_i\ell_i(X))_i\bigr).
\]
It commutes with the mapping class action (and with a simultaneous permutation of end labels). On the positive-length part its group of units acts freely and transitively on each capping fibre. The zero vector gives the canonical capping retraction. The maps $X\mapsto\Phi^{-1}(\operatorname{cap}(X),(1-t)\ell(X))$ give a strong deformation retraction onto the punctured locus.
\end{corollary}
\begin{proof}
All assertions follow by conjugating coordinatewise multiplication and straight contraction by the specified homeomorphism $\Phi$. Its definition uses only canonical capping and the intrinsic hyperbolic boundary lengths. It therefore makes no choice of a pants decomposition. It need not agree with the operation that holds unrelated Fenchel--Nielsen coordinates fixed while multiplying selected lengths.
\end{proof}
